The organization \must{} maintain an accurate inventory of software used on organizational systems or to access organizational data.

The organization \should{} maintain an approved software list identifying software that is authorized for general use, restricted use, or use only on specifically approved systems.

Software \should{} be assigned a status such as approved, approved with restrictions, pending review, prohibited, unsupported, or retired.

The software inventory \must{} include, where applicable:
\begin{itemize}
    \item Operating systems and system components
    \item Applications and productivity software
    \item Security, backup, monitoring, and management tools
    \item Databases, runtimes, libraries, frameworks, and dependencies
    \item Firmware and embedded software
    \item Scripts, automation tools, and internally developed software
    \item Software-as-a-service applications and cloud services
    \item Mobile applications used to access organizational systems or data
    \item Container images and packages
    \item Browser extensions and add-ons
\end{itemize}

Each software record \must{}, to the extent applicable, include:
\begin{itemize}
    \item Software name and purpose
    \item Publisher, developer, or service provider
    \item Version and edition
    \item Installation location or associated system
    \item Software owner and technical administrator
    \item License type, license quantity, and renewal date
    \item Data classification and types of data handled
    \item Business criticality
    \item Dependencies and integrations
    \item Support status and end-of-support date
    \item Approved source or distribution method
    \item Date of last review
\end{itemize}

Software \must{} be approved before it is installed, deployed, or used to process organizational data. Approval \must{} consider the software's purpose, security features, licensing requirements, privacy and data-handling practices, compatibility, support status, and operational risk.

Users \mustnot{} install or use unauthorized software on organizational systems. Users \mustnot{} use software obtained from untrusted or unofficial sources unless it has been reviewed and approved by an authorized administrator.

Software installation \must{} be limited to authorized personnel or approved management systems. Administrative privileges \should{} not be granted to users solely to permit software installation.

The organization \must{} maintain a process for identifying unsupported, obsolete, vulnerable, or end-of-life software. Unsupported software \must{} be upgraded, replaced, isolated, or removed. If immediate remediation is not possible, the system owner \must{} document and approve an exception that identifies the risk, compensating controls, responsible owner, and expiration or review date.

Software \must{} be kept reasonably current through approved patches, updates, and upgrades. Security updates \must{} be prioritized according to the severity of the vulnerability, exposure of the affected system, availability of mitigations, and business criticality.

Software updates \should{} be tested before broad deployment when testing is practical and the risk of incompatibility is significant. Emergency security updates may be deployed without standard testing when delaying deployment would create a greater risk.

The organization \should{} use centralized software deployment, endpoint management, package management, or equivalent technical controls to:
\begin{itemize}
    \item Identify installed software
    \item Detect unauthorized software
    \item Track software versions
    \item Deploy approved updates
    \item Remove unsupported or unauthorized software
    \item Report systems that do not meet software requirements
\end{itemize}

Open-source software, libraries, packages, and external dependencies \must{} be obtained from approved sources and maintained at supported versions where feasible. Internally developed software \must{} identify its significant external dependencies and \should{} be reviewed for known vulnerabilities before release.

Software used to handle confidential, regulated, or otherwise sensitive information \must{} be evaluated for appropriate access controls, encryption, logging, retention, and data-sharing behavior.

Software licenses \must{} be managed in accordance with applicable licensing terms. The organization \must{} maintain sufficient records to demonstrate that software use is authorized and that license limits are not exceeded.

Software-as-a-service applications \must{} have an identified business owner and technical or administrative contact. Before use, the organization \should{} evaluate the service's authentication, access control, data protection, backup, logging, incident notification, termination, and data-return capabilities.

Software inventory records \must{} be reviewed:
\begin{itemize}
    \item At least annually
    \item Before and after major software deployments
    \item When software is replaced, upgraded, or retired
    \item When a vendor ends support for a product
    \item After a significant vulnerability or security incident
    \item When licensing, contractual, or regulatory requirements change
\end{itemize}

Software that is no longer approved, supported, required, or licensed \must{} be removed or formally excepted. Associated accounts, integrations, credentials, scheduled tasks, service connections, and stored data \must{} be reviewed and disabled or removed as appropriate.

Retired software records \should{} include:
\begin{itemize}
    \item Software name and version
    \item Systems or users affected
    \item Retirement date
    \item Reason for retirement
    \item Replacement software, if applicable
    \item Data migration or disposal actions
    \item Removal or decommissioning approval
\end{itemize}

Changes to critical software installations, configurations, integrations, and versions \should{} be documented. Significant changes \must{} follow the organization's change-management procedures.